\section{Limitations and ongoing extensions}
\label{sec:limitations}

\subsection{What the current system cannot express}

The current §6-structural emitter handles cyclic proofs whose recursion
fits Lean's structural-recursion check. Several classes of proofs lie
outside this scope.

\paragraph{Non-structural recursion.}
The most visible limitation is exposed by merge-sort. The auxiliary lemma
\(\code{length (merge xs ys) = length xs + length ys}\) recurses on the
result of a Decidable comparison: when \(x \le y\), it descends on
\code{xs}; otherwise on \code{ys}. Neither argument is structurally smaller
across both cases, so Lean's structural-recursion check rejects the
underlying \code{def} produced by Phase~B, and the structural §6 emitter
cannot witness termination on any single variable. This is precisely the
case the dispatcher (§\ref{sec:dispatch}) detects and routes to the
well-founded backend, where a lexicographic measure
\((\code{xs.length}, \code{ys.length})\) takes over; merge-sort additionally
exposes a residual arithmetic-discharge gap (the \code{+1} step the hand
proof leaves to \code{omega}) that is independent of the dispatch.

\paragraph{Decidable case-splits.}
The original \code{cyc\_cases} tactic recorded only inductive-type case
splits, so a proof body containing \code{by\_cases h : P} with back-edges
in each branch was opaque to the event recorder. The \code{cyc\_by\_cases}
tactic (\file{CyclicTactic/Tactic.lean}) now closes this gap: it records a
Decidable split as a \code{.dCaseSplitStart} event the same way
\code{cyc\_cases} records inductive ones, so the back-edges within each
branch enter the size-change check and the proof can dispatch to WF
emission (§\ref{sec:dispatch}).

\paragraph{DAG-shaped cyclic proofs.}
The system assumes back-edges target ancestors only. Cyclic proofs that
share intermediate sub-derivations across multiple branches (DAG-shaped
proofs in the sense of Brotherston's PhD~\cite{brotherston2006phd}) cannot
be expressed directly. The standard workaround---factoring shared
sub-derivations into separate \code{cyclic\_thm}s or \code{cyclic\_mutual}
entries---works for most practical proofs but is not always idiomatic. We
discuss the theoretical equivalence in §\ref{sec:dag-discussion} below.

\paragraph{Multi-companion in a single theorem.}
The §6 data model assumes companion = root: every back-edge in a single
\code{cyclic\_thm} body cycles to the same companion (the root sequent).
Proofs whose natural shape uses multiple \code{cyclic} declarations within
one theorem body (with different buds targeting different companions) are
not expressible. \code{cyclic\_mutual} offers a workaround by splitting
into mutually-recursive entries, but the surface ergonomics differ.

\paragraph{Within-entry cycles in \code{cyclic\_mutual}.}
The mutual emitter currently skips §6 augmentation, because
\code{annotateTree}'s Lemma 5.9 unfolding does not terminate on
cross-theorem buds. A \code{cyclic\_mutual} block whose entries contain
within-entry back-edges (in addition to cross-companion edges) is therefore
not supported. The §6 framework's extension to mutual systems is, to our
knowledge, an open theoretical question.

\paragraph{Proposition 5.8 verification.}
Our system uses Wehr~3.2.4's lex-order finder (Theorem~3.2.4 of
\cite{wehr2025cyclic}) as a check on the induction order, complementary to
the Grotenhuis-Otten Theorem~5.2 progressing-name verification. The paper's
Proposition~5.8 gives a stronger verification of the induction order that
we do not implement; we delegate this to Lean's structural-recursion check.
Implementing Proposition~5.8 directly would turn the current diagnostic
PASS/FAIL into a real proof obligation.

\paragraph{Multi-sort extension (§8).}
Grotenhuis-Otten §8 extends the unravelling algorithm to multi-sorted
inductive systems. We currently handle single-sorted cyclic proofs;
extending to §8 would let cyclic proofs span fundamentally distinct
inductive types as their recursion structure.

\subsection{Tree-shape versus DAG-shape: a discussion}
\label{sec:dag-discussion}

It is worth being precise about what is lost when we restrict to
tree-shaped cyclic proofs with back-edges to ancestors only. Sprenger and
Dam's Theorem~5 and its descendants establish that tree-shaped and
DAG-shaped cyclic proofs prove the same sequents: any DAG can be unfolded
to a tree by duplicating shared sub-derivations, with potentially
exponential blow-up. So provability is not affected.

What \emph{is} affected:

\begin{itemize}
  \item \textbf{Proof size.} A DAG with many shared sub-derivations
    unfolds to a tree exponentially larger.
  \item \textbf{Sharing-natural patterns.} Symmetry-driven proofs,
    cross-recursive lemmas with shared support, and confluence-style
    arguments naturally have DAG shape.
  \item \textbf{Lateral cycles.} A back-edge from one branch of a
    proof tree to a node in another branch (not an ancestor) cannot be
    represented.
\end{itemize}

For the Lean tactic-mode setting we target, this restriction is less
constraining than it sounds. Lean idiom favours factoring shared
sub-derivations into separate \code{theorem}s or local \code{have}
binders---both of which absorb the sharing-naturalness patterns above. The
remaining genuinely lateral case (a single proof tree with a non-ancestor
back-edge) is rare in practice and can typically be expressed by promoting
the lateral target to a separate \code{cyclic\_thm}.

This trade-off is most visible against Cyclist's separation-logic
benchmark (\texttt{benchmarks/sl/}), where DAG shape is heavily used by the
proof-search procedure and our system would struggle. For Cyclist's
first-order benchmark and for typical Lean proof work, tree-shape is a
defensible architectural choice rather than a deficiency.

\subsection{The dispatcher: structural, well-founded, and Löb emission}
\label{sec:dispatch}

The Grotenhuis-Otten §6 algorithm assumes structural emission is always
possible. It is not: merge-sort and the swap-recursion of Cyclist's
Test~14 (§\ref{sec:limitations}) admit no single-variable structural
order. The system therefore \emph{dispatches} among emission strategies on
the basis of the size-change data, rather than committing to structural
emission unconditionally. Two of the three strategies are now implemented
on the \code{wf} branch; the third remains under separate development. We
describe each, then state the characterisation theorem they converge on.

\paragraph{The dispatch predicate.}
The choice between structural and well-founded emission is made by a
single Boolean predicate, \code{canStructural}
(\file{CyclicTactic/Tactic.lean}):

\begin{lstlisting}
def canStructural (tree : ProofTree)
    (labeled : List (String × SCGraph)) (arity : Nat) : Bool :=
  match InductionOrder.findInductionOrder tree labeled arity with
  | none     => false
  | some asg => (InductionOrder.lexOrderFromAssignment asg).length <= 1
\end{lstlisting}

\code{findInductionOrder} runs the Wehr~3.2.4 lex-order finder over the
bud-companion strongly-connected components; it returns an assignment of
progressing argument positions, or \code{none} when no consistent order
exists. \code{canStructural} returns \code{true} exactly when that
assignment, after deduplication, uses \emph{at most one} argument
position---i.e.\ when a single \code{induction} on one variable carries the
whole recursion, which is precisely the shape Theorem~6.1's structural
emitter (§\ref{sec:impl-emission}) produces. When the order needs two or
more positions (a genuine lexicographic descent), or does not exist at all,
the predicate returns \code{false} and the proof is routed to the
well-founded backend. Merge-sort's \code{merge} fails the predicate
because its \(x\le y\) arm descends on \code{xs} (position~0) and its
\(x>y\) arm on \code{ys} (position~1): no single position is
preserved-and-progressing across both arms, so \code{findInductionOrder}
returns \code{none}.

\paragraph{Well-founded (WF) emission.}
When \code{canStructural} is \code{false}, the system emits
\code{def \dots\ termination\_by <measure>} with a
\code{decreasing\_by simp\_wf; omega} discharge, via
\code{WFEmit.translate} (\file{CyclicTactic/WFEmit.lean}). The measure is
synthesised from the per-back-edge size-change graphs by
\code{synthMeasure} (\file{CyclicTactic/Measure.lean}), which tries four
schemas in increasing generality (the lex-ranking ideas are
Lee~2009's~\cite{lee2009ranking}):

\begin{enumerate}
  \item \textbf{Lex over a permutation} of all argument positions---a
    strict self-loop at some position, nonstrict at all earlier ones
    (classic lex termination, e.g.\ Ackermann).
  \item \textbf{Lex over an ordered subset}, allowing non-participating
    positions to be dropped from the tuple.
  \item \textbf{Sum measure} \(a_0+\dots+a_{n-1}\), for swap-style
    recursions where the bag of arguments shrinks as a whole.
  \item \textbf{Closure-witness greedy lex}: compute the SCT composition
    closure, take each idempotent's strict-self-loop positions as a
    soundness witness, and greedily build a lex order covering the most
    uncovered idempotents at each step. This is a flat-arity
    specialisation of the Grotenhuis-Otten Definition~5.1 reset-proof
    construction.
\end{enumerate}

Every schema validates its candidate against the \emph{input} graphs, not
merely the closure idempotents, because Lean's \code{termination\_by}
requires every recursive call to decrease, not just every cycle. The
result is a \code{Measure}, which is either \code{lex order} or
\code{sum arity}. The Decidable case-splits that such proofs require
(merge-sort's \code{x \(\le\) y}) are now recorded by the \code{cyc\_by\_cases}
tactic, which emits a \code{.dCaseSplitStart} event exactly as
\code{cyc\_cases} does for inductive splits; this closes the gap noted
in §\ref{sec:limitations} above. The swap recursion of Test~14
(\code{N2(x,y)} with recursive call \code{N2(y,x)}) goes through
end-to-end: \code{synthMeasure} produces \code{sum 2}, rendered
\code{termination\_by x + y}, and the emitted \code{n2\_holds} is a
kernel-checked term (\file{CyclicTactic/Examples/WFSuite.lean}).

\paragraph{Löb-induction emission.}
A theoretically stronger alternative replaces structural and WF emission
entirely with \emph{Löb induction} via a step-indexed
embedding~\cite{atkey2013productive, jung2015iris}. Every SCT-validated
cyclic proof can be emitted as a Löb application, using a \(\Later\)
modality defined by \(\Later P\, n := \forall m < n, P\, m\) and the Löb
axiom \((\forall n.\, \Later P\, n \to P\, n) \to \forall n.\, P\, n\)
proved from \code{Nat.strongRec}. The advantage is universality: Löb
emission applies to any SCT-valid cyclic proof, with no case analysis on
which strategy fits. The disadvantage is artefact form: emitted proofs are
about \(\Later P\) rather than \(P\) directly, which is less idiomatic for
Lean users. This extension is under active development on a separate
\code{lob} branch in a sibling repository, and is the only one of the
three not yet integrated into the main dispatcher.

\paragraph{The characterisation theorem.}
The dispatcher \code{canStructural} and the synthesiser \code{synthMeasure}
are total, decidable functions on the §6 size-change data; what is
\emph{conjectural} is that their decisions agree with Lean's
own acceptance of the emitted artefact. We separate the two.

\begin{conjecture}[Dispatch soundness]
\label{conj:dispatch}
Let \(T\) be an SCT-valid \code{ProofTree} with labelled size-change
graphs \(G\) and arity \(n\), and write
\(A=\code{computeAug}\,T\) for its §6 augmentation.
\begin{enumerate}
  \item If \(\code{canStructural}\,T\,G\,n = \mathtt{true}\), then the
    structural emitter \(\code{Emit.translate}\,T\) produces a term Lean's
    structural-recursion check accepts.
  \item If \(\code{canStructural}\,T\,G\,n = \mathtt{false}\) and
    \(\code{synthMeasure}\,G\,n = \code{some}\,M\), then the well-founded
    emitter \(\code{WFEmit.translate}\,M\,T\) produces a \code{def} whose
    \(\code{termination\_by}\,M\) obligation \code{decreasing\_by} discharges.
  \item The residual case
    \(\code{canStructural} = \mathtt{false}\) and
    \(\code{synthMeasure} = \code{none}\)---an SCT-valid proof whose
    descent is captured by no lex-or-sum measure---is exactly the regime
    in which the first two strategies fail and Löb emission is required.
\end{enumerate}
\end{conjecture}

\begin{conjecture}[Löb universality]
\label{conj:lob}
For every SCT-valid \code{ProofTree} \(T\), the Löb emitter
\(\code{Emit.translateLob}\,T\) produces a Lean-acceptable term.
\end{conjecture}

The two extensions thus carve up the SCT-valid proofs into a structural
core (Conjecture~\ref{conj:dispatch}.1), a well-founded middle
(Conjecture~\ref{conj:dispatch}.2), and a Löb-only remainder
(Conjectures~\ref{conj:dispatch}.3 and~\ref{conj:lob}). Clauses~1 and~2
are now backed by working emitters and an example suite; proving that the
\emph{syntactic} predicates \code{canStructural} and \code{synthMeasure}
soundly approximate Lean's \emph{semantic} acceptance---and that the Löb
remainder is genuinely universal---are the load-bearing claims that anchor
the next stage of the work. A proof of Conjecture~\ref{conj:dispatch}.1
would proceed by showing the Wehr~3.2.4 single-position assignment is
exactly the witness Lean's structural-recursion elaborator searches for;
of clause~2, by appeal to the soundness of the Lee~2009 ranking schemas
against the input graphs; and of Conjecture~\ref{conj:lob}, by exhibiting
the step-indexed model in which \(\Later\) is the strict-order
approximation and reading SCT validity as the descent that makes the Löb
premise discharge.

\subsection{Other engineering follow-ups}

Several smaller items are not algorithmically interesting but are
worth listing:

\begin{itemize}
  \item The \code{CyclicTactic/Examples/CyclistComparison.lean} file
    has a pre-existing stack overflow in
    \code{SubjectTerm.toString}, triggered by one of the later examples
    in the file. The crash predates the §6 work and likely involves
    infinite term traversal when an inductive-predicate hypothesis is
    recorded. Investigating this is independent of the §6 / WF / Löb
    work but blocks bare \code{lake build}.

  \item The system does not yet support \code{cyclic\_thm} bodies whose
    structure is \emph{not} a sequence of \code{cyc\_cases} \(+\)
    \code{back} \(+\) leaves. In particular, intermediate \code{cyclic}
    re-declarations within a single body are not handled, and
    interleaving \code{cyclic\_thm}-style tactics with other custom
    tactics from libraries is untested.

  \item Diagnostic output is verbose but ad hoc; a structured
    \code{[cyclic\_thm]} info schema (so external tooling can parse the
    dispatcher's decisions) would make integration with editor extensions
    easier.

  \item Performance of the §6 augmentation has not been profiled. For
    very deep proof trees the \code{annotateTree} \(+\)
    \code{unfoldFixpoint} cost could become significant; we have not
    observed this on the example suite.
\end{itemize}
